\documentclass[11pt]{article}
\usepackage[margin=0.9in]{geometry}
\usepackage{amsmath,amssymb,amsthm,booktabs,needspace}
\usepackage[hidelinks]{hyperref}
\setlength{\emergencystretch}{2em}
\newtheorem{theorem}{Theorem}
\newtheorem{lemma}[theorem]{Lemma}
\newcommand{\rank}{\operatorname{rank}}
\title{Compact presentations and exact four-port assemblies for two double covers}
\author{Internal research working proof}
\date{17 November 2026}
\begin{document}
\maketitle

The preceding mutation theorem reduces eight retained ordinary cover-rank
questions to two. This note supplies short, independently replayed
presentations for those groups and identifies a common pair of formal
fragment groups whose four port identifications reconstruct each source.
The problem is now concentrated in those identifications. Neither rank
two nor rank three is proved; no new knot equality follows. The uniform
algebra below is a written working proof, unformalized and externally
unreviewed. No worldwide novelty claim is made.

\section{The two short source presentations}

Write $G_A=\pi_1(\Sigma_2(13a_{129}))$ and
$G_B=\pi_1(\Sigma_2(13a_{135}))$. The retained mutation theorem transfers
their ordinary ranks to the respective blocks
$\{129,136,194,197\}$ and $\{135,137,196,198\}$, where every name has
prefix $13a$. Each member retains its own four-bridge upper certificate.

\begin{theorem}\label{thm:short}
The source groups have the following exact presentations:
\begin{align*}
G_A=\langle a,b,c\mid{}&
c^{-2}b^{-1}c^{-2}b^{-1}c^{-2}b,\\[-2pt]
&c^{-1}a^{-2}b^{-2}c^{-1}a^{-5}b^{-2},\\[-2pt]
&c^{-1}a^3b^{-2}ca^3b^{-2}cb^{-2}\rangle;\\[4pt]
G_B=\langle a,b,c\mid{}&
c^{-3}b^{-3}a^{-1}c^{-3}a^2,\\[-2pt]
&c^{-2}a^3b^3cab^3a^2,\\[-2pt]
&b^{-5}a^2b^{-2}a^2\rangle.
\end{align*}
Their relator lengths are respectively $(9,13,15)$ and $(12,15,11)$.
Their abelianizations are $C_{391}$ and $C_{401}$, and each maps onto
the retained nonabelian group of order $5616$. Consequently
$2\le\rank G_A,\rank G_B\le3$; the actual ranks remain undetermined.
\end{theorem}

The input is the literal pinned thirteen-crossing PD code for each
source. With legacy arc indices $0,\ldots,12$, the normalized core
presentation has $x_0=1$ and one relation
\[
 x_o x_a^{-1}x_o x_b^{-1}=1
\]
at each crossing with overarc $o$ and underarcs $a,b$. Its identification
with the double-cover group is the previously retained classical
core-group input, credited through Silver--Traldi--Williams. It is
not inferred from the new simplification.

Every accepted move is recorded: solving a relator for a generator
occurring once; applying a free automorphism with an explicit inverse;
or replacing a cyclic relator subword $U$ by $V^{-1}$ using another
relator $UV$. The last replacement is precisely multiplication by
$(UV)^{-1}$ after choosing the displayed cyclic representatives.
Free reduction, cyclic conjugation, inversion, and removal of trivial
or duplicate relators preserve the normal closure. These are exact
Tietze and free-basis changes. A strictly decreasing length search
chooses moves but supplies no rank lower bound when it stops.

For $G_A$ the trace has nine eliminations, five automorphisms and three
relator replacements. For $G_B$ the counts are nine, four and six.
\texttt{verify\_compact\_covers.py} independently reconstructs the core
words, replays every move, and checks both recorded generator maps.
The surviving labels $(4,6,9)$ and $(4,9,12)$ are renamed $(a,b,c)$
in Theorem~\ref{thm:short}; these are changed free-basis generators,
not assertions that the final generators equal the original arcs.
Integral determinantal divisors give Smith factors $(1,1,391)$ and
$(1,1,401)$. Transported projective permutations satisfy every original
and final relation and generate groups of order $5616$.
These checks prove the finite identities used in the theorem.

\section{A common pair of formal fragment groups}

For each root, split its source PD into the first seven and last six
crossings. Each part has four boundary edge labels, exactly
$12,15,17,26$. Identifying only overpassing edge pairs gives nine arc
classes on the inner fragment and eight on the outer fragment. Normalize
the first class in each part to one and impose its core crossing
relations. These formal groups have independently replayed presentations
\begin{align*}
 P&=\langle u,v\mid v^{-2}u^2v^{-1}u^{-3}v^{-1}u^2\rangle,\\
 Q&=\langle x,y\mid y^{-2}x^{-2}y^{-2}x\rangle.
\end{align*}
Both roots give exactly these two abstract presentations after the
recorded renamings. Each group has abelianization $\mathbb Z$.
A retained cyclic-order-reversing ribbon map identifies the two outer
fragments with their ports tracked. It is not an orientation-preserving
diagram identification or an independently certified torus decomposition.

Define $w=v u^{-2}v$. The normalized port words are:
\begin{center}
\begin{tabular}{llll}
\toprule
Edge & Inner $p_e$ & Outer $q_e$ for $A$ & Outer $q_e$ for $B$\\
\midrule
$12$ & $w^3u$ & $1$ & $1$\\
$15$ & $u^2v^{-1}u^3$ & $y^2$ & $y^2$\\
$17$ & $vu$ & $x^{-3}y$ & $x^3y$\\
$26$ & $1$ & $y^{-1}xy^{-2}x$ & $y^{-1}x^{-2}y^{-2}x^{-2}$\\
\bottomrule
\end{tabular}
\end{center}
All words are in the fixed bases of $P,Q$ above. The internal labels
used to obtain those bases differ between sources; the retained
generator maps account for this.

\begin{theorem}[Exact algebraic assembly]\label{thm:assembly}
For $X=A,B$, with the corresponding outer column in the table,
\[
 G_X\cong
 (P*Q*\langle z\rangle)\big/
 \langle\!\langle p_e zq_e^{-1}:e=12,15,17,26\rangle\!\rangle.
\]
This is an exact group-presentation identity. No injectivity of a
boundary subgroup or topological JSJ assertion is included.
\end{theorem}

\begin{proof}
First undo the recorded fragment reductions. Write the inner normalized
arc variables as $u_i$, with $u_1=1$, and the outer ones as $v_j$, with
$v_1=1$. Introduce $z$ and set the unnormalized inner variables to
$s_i=u_i z$. This is a free-basis change: $s_1=z$ and
$u_i=s_i s_1^{-1}$. At every inner crossing, the right factors cancel:
\[
 (u_o z)(u_a z)^{-1}(u_o z)(u_b z)^{-1}
 =u_o u_a^{-1}u_o u_b^{-1}.
\]
The four port relations now simply identify the incident inner $s_i$
and outer $v_j$ variables. Independent edge-path reconstruction checks
that the seventeen fragment arc classes, joined by these four
independent identifications, are exactly the thirteen complete source
arc classes. Thus the remaining thirteen crossing relators are the
full source core relations, normalized at the outer base arc.

The outer base is source arc $5$, rather than arc $0$ used earlier.
Changing the normalized base from $j$ to $k$ sends each arc variable
$x_i$ to $y_i y_j^{-1}$, where $y_k=1$. The inverse sends $y_i$ to
$x_i x_k^{-1}$. Their compositions freely reduce to the identity, and
the same displayed cancellation checks every crossing relation.
Therefore the two normalizations give isomorphic source groups.
Reapplying the independently replayed compact fragment moves gives
the displayed assembly.
\end{proof}

The assembly initially has five generators and six relators. Its own
recorded reductions return three-generator presentations with cyclic
abelianizations $391$ and $401$. The direct presentations in
Theorem~\ref{thm:short} and the port assemblies are two exact descriptions
of each source; neither asserts minimality.

\section{Independent checks and the boundary of the result}

The compact verifier imports no producer simplification routine.
The assembly verifier reuses that independent word replay and builds
fragment arc classes by graph traversal, rather than the producer's
union construction. It checks all four literal ports and their words,
all fragment and assembly traces, and the seventeen-to-thirteen arc
identification. A legacy PD parser is shared only for old arc names,
with geometric paths checked; this dependence is disclosed.

Three known algebraic controls check ranks two and three
through exact simplification and $\mathbb F_3$ homology: two and three
free factors $C_3$, and a rebased two-factor presentation with an extra
eliminable generator. Five compact corruptions and four additional
port/assembly corruptions are rejected, including a false assertion
that all source arcs have been generated from the outer fragment.

Indeed, initializing the known outer arc words and using only
$x_b=x_o x_a^{-1}x_o$ whenever the overarc and one underarc are known
stalls on source arcs $0,1,2,3,4$ in both roots. That verifies the limit
of this particular deduction. It proves neither rank three nor absence
of another two-generator certificate. Two-generated pieces alone do
not make the assembled group two-generated.

The selected abstract Tait graphs also admit exact series/parallel
reductions to $K_4$. One branch contains seven original crossings,
another adjacent branch two, and the other four one each. The seven-edge
network has tree/separating-forest counts $(19,14)$ for $A$ and $(14,19)$
for $B$. Enumerating the sixteen $K_4$ spanning trees gives the whole
tree-count formula $11T+13F$, hence $391$ and $401$. These are graph
network data. Their ratios have not been certified as rational-tangle
slopes; series/parallel arithmetic alone does not establish the needed
geometric filling model.

Small, separately frozen software probes give only leads. Candidate
cyclic-double-cover fillings have the expected homology, but numerical
shapes are degenerate and a single Regina standard-recognition pass
recognizes neither target. The figure-eight control is recognized as
$L(5,2)$. Neither target's peripheral filling provenance or exact
core-to-triangulation transport is certified by those probes. No
hyperbolicity, graph-manifold, Seifert or JSJ conclusion is drawn.

The next rank question is to determine the boundary peripheral words
and the nature of the two pieces, then certify the relevant gluing
and injectivity before using a manifold or group-amalgamation rank
theorem. Brunner/Ito presentations and Boileau--Weidmann rank papers
are source leads, not imported classifications in this result.
The finite packages are \texttt{compact\_cover\_verification.json} and
\texttt{port\_assembly\_verification.json}. The catalogue remains at
42 sources. All original sixteen target equalities, full MRC and the
broader source and priority audits remain open.

\end{document}
